\section*{Chapitre \ref{ch:b2:randomvar} --- Variables aléatoires discrètes}

\begin{solution}{exo:b2:randomvar:1}
\emph{Binomiale~:} $X = \sum_{i=1}^n X_i$ avec des $X_i$ de Bernoulli
indépendantes~; $V(X_i) = \E(X_i^2) - \E(X_i)^2 = p - p^2$, et les \omterm{def:b2:randomvar:variance}{variances}
de variables indépendantes s'ajoutent
(\cref{thm:b2:randomvar:variancerules})~:
\[
\E(X) = np, \qquad V(X) = np(1-p) .
\]
\emph{Poisson~:} $\E\bigl(X(X-1)\bigr) =
\sum_{k\geq2}k(k-1)e^{-\lambda}\frac{\lambda^k}{k!} = \lambda^2
e^{-\lambda}\sum_{j\geq0}\frac{\lambda^j}{j!} = \lambda^2$, donc
\[
V(X) = \E(X^2) - \E(X)^2
= \lambda^2 + \lambda - \lambda^2 = \lambda .
\]
\emph{Géométrique} ($q = 1 - p$)~: en dérivant deux fois
$\sum_{k\geq0}q^k = \frac{1}{1-q}$ à l'intérieur du disque
(\cref{ch:b2:powerseries}), $\sum_{k\geq2}k(k-1)q^{k-2} =
\frac{2}{(1-q)^3}$, donc
\[
\E\bigl(X(X-1)\bigr) = pq\sum_{k\geq2}k(k-1)q^{k-2}
= \frac{2q}{p^2},
\qquad
V(X) = \frac{2q}{p^2} + \frac1p - \frac{1}{p^2}
= \frac{q}{p^2} = \frac{1-p}{p^2} .
\]
\end{solution}

\begin{solution}{exo:b2:randomvar:2}
\emph{Somme~:} pour $n \in \N$, par disjonction et indépendance,
\[
\P(X + Y = n)
= \sum_{k=0}^n \P(X = k)\P(Y = n - k)
= e^{-(\lambda + \mu)}\frac{1}{n!}
\sum_{k=0}^n \binom nk \lambda^k\mu^{n-k}
= e^{-(\lambda+\mu)}\frac{(\lambda + \mu)^n}{n!}
\]
par la formule du binôme~: $X + Y \sim \mathcal{P}(\lambda + \mu)$.
\emph{\omterm{def:b2:randomvar:law}{Loi} conditionnelle~:} pour $0 \leq k \leq n$,
\[
\P(X = k \mid X + Y = n)
= \frac{\P(X = k)\P(Y = n - k)}{\P(X + Y = n)}
= \binom nk
\Bigl(\frac{\lambda}{\lambda+\mu}\Bigr)^{k}
\Bigl(\frac{\mu}{\lambda+\mu}\Bigr)^{n-k} ,
\]
la \omterm{def:b2:randomvar:law}{loi} binomiale $\mathcal{B}\bigl(n,
\frac{\lambda}{\lambda+\mu}\bigr)$~: étant donné le comptage total, chaque
\omterm{def:b2:proba:space}{événement} «~choisit~» indépendamment la première source avec probabilité
proportionnelle à son taux.
\end{solution}

\begin{solution}{exo:b2:randomvar:3}
Lorsque $k$ jouets manquent encore, chaque nouvelle boîte apporte un nouveau jouet
avec probabilité $\frac kn$, indépendamment du passé~: le
temps d'attente $W_k$ du prochain nouveau jouet est géométrique
$\mathcal{G}\bigl(\frac kn\bigr)$, avec $\E(W_k) = \frac nk$, et
$T_n = W_n + W_{n-1} + \dots + W_1$ (la première boîte donne toujours un nouveau
jouet~: $W_n = 1$, cohérent avec $\E = n/n$). Par linéarité,
\[
\E(T_n) = \sum_{k=1}^n \frac nk = n\sum_{k=1}^n\frac1k
\sim n\ln n ,
\]
en utilisant $\sum_{k\leq n}\frac1k = \ln n + \gamma + o(1)$
(\cref{ch:b2:comparison}). Collectionner les derniers jouets est ce qui
coûte~: la moitié des boîtes va à la dernière poignée.
\end{solution}

\begin{solution}{exo:b2:randomvar:4}
Ponctuellement, $X(\omega) = \#\{n \geq 1 : X(\omega) \geq n\} =
\sum_{n\geq1}\mathbf{1}_{X \geq n}(\omega)$. La famille double
$\bigl(\mathbf{1}_{X \geq n}(\omega)\,\P(\{\omega\})\bigr)_{n,
\omega}$ est positive, donc Fubini pour les familles
(\cref{ch:b2:series}) s'applique inconditionnellement~: en sommant d'abord en
$n$ on obtient $\E(X)$, en sommant d'abord en $\omega$ on obtient $\sum_n \P(X
\geq n)$~; les deux sont simultanément finies et égales. Pour $X \sim
\mathcal{G}(p)$~: $\P(X \geq n) = q^{n-1}$ ($q = 1-p$), donc $\E(X) =
\sum_{n\geq1}q^{n-1} = \frac{1}{1 - q} = \frac1p$.
\end{solution}

\begin{solution}{exo:b2:randomvar:5}
\emph{Symétrie~:} la $i$-ème boule tirée est une boule \omterm{def:b2:funcseq:def}{uniformément}
aléatoire de l'urne (chacune des $N$ boules a la même probabilité d'atterrir en
position $i$ de l'\omterm{def:b2:structures:generated}{ordre} de tirage), donc $\P(Y_i = 1) = \frac MN =
p$ et $\E(X) = np$ par linéarité --- aucune indépendance nécessaire.

\emph{\omterm{def:b2:randomvar:variance}{Covariance}~:} pour $i \neq j$, $\E(Y_iY_j) = \P(\text{les tirages }
i, j \text{ sont tous deux blancs}) = \frac{M(M-1)}{N(N-1)}$ (les couples
ordonnés de positions distinctes reçoivent un couple ordonné de boules distinctes,
\omterm{def:b2:funcseq:def}{uniformément}). Donc
\[
\operatorname{Cov}(Y_i, Y_j)
= \frac{M(M-1)}{N(N-1)} - \frac{M^2}{N^2}
= \frac{M(N - M)}{N^2}\cdot\frac{-1}{N-1}
= -\frac{p(1-p)}{N-1} < 0 :
\]
tirer une boule blanche raréfie les blanches pour les autres tirages.
Par le~\cref{thm:b2:randomvar:variancerules},
\[
V(X) = np(1-p) + n(n-1)\Bigl(-\frac{p(1-p)}{N-1}\Bigr)
= np(1-p)\,\frac{N - n}{N - 1} \leq np(1-p) :
\]
l'échantillonnage sans remise a la même moyenne mais une \omterm{def:b2:randomvar:variance}{variance} \emph{plus petite}
qu'avec remise (égalité seulement pour $n = 1$), les corrélations
négatives agissant comme un stabilisateur. Pour $n = N$ la
\omterm{def:b2:randomvar:variance}{variance} s'annule~: le comptage est alors déterministe.
\end{solution}

\begin{solution}{exo:b2:randomvar:6}
En développant autour de $m = \E(X)$~:
\[
\E\bigl((X - c)^2\bigr)
= \E\bigl((X - m)^2\bigr) + 2(m - c)\,\E(X - m) + (m - c)^2
= V(X) + (m - c)^2 ,
\]
minimale exactement en $c = m$ de valeur $V(X)$ --- l'\omterm{def:b2:randomvar:expectation}{espérance} est
le meilleur prédicteur constant en moyenne quadratique.

Si $\P(X = m) = 1$ alors $(X - m)^2$ s'annule avec probabilité
$1$, donc $V(X) = 0$ (la famille définissante a des termes nuls sauf sur un
ensemble négligeable). Réciproquement, si $V(X) = 0$, Tchebychev
(\cref{thm:b2:randomvar:markov}) donne $\P\bigl(\abs{X - m} \geq
\frac1n\bigr) \leq n^2\,V(X) = 0$ pour tout $n$~; les \omterm{def:b2:proba:space}{événements}
$\bigl\{\abs{X - m} \geq \frac1n\bigr\}$ croissent vers $\{X \neq
m\}$, donc la continuité monotone (\cref{thm:b2:proba:continuity})
donne $\P(X \neq m) = 0$.
\end{solution}

\begin{solution}{exo:b2:randomvar:7}
$\E(S_n) = \frac n2$ et $V(S_n) = \frac n4$.
\emph{Markov~:} $\P\bigl(S_n \geq \frac{3n}4\bigr) \leq
\frac{n/2}{3n/4} = \frac23$ --- une borne constante, inutile pour
$n$ grand.
\emph{Tchebychev~:} l'\omterm{def:b2:proba:space}{événement} implique $\abs{S_n - \frac n2} \geq
\frac n4$, donc la probabilité est $\leq \frac{n/4}{(n/4)^2} =
\frac4n$ --- décroît, mais seulement polynomialement.
\emph{Chernoff~:} par indépendance, $\E(e^{tS_n}) =
\prod_{i=1}^n\E(e^{tX_i}) = \bigl(\frac{1 + e^t}{2}\bigr)^n$, et
Markov appliquée à $e^{tS_n} \geq e^{3nt/4}$ donne, pour tout $t >
0$,
\[
\P\Bigl(S_n \geq \frac{3n}4\Bigr)
\leq \Bigl(\frac{1 + e^t}{2}\Bigr)^n e^{-3nt/4}
= \exp\Bigl(n\bigl(\ln\tfrac{1 + e^t}{2} - \tfrac{3t}4\bigr)\Bigr).
\]
Minimisons l'exposant~: $\frac{\dd}{\dd t}\ln\frac{1+e^t}{2} =
\frac{e^t}{1 + e^t} = \frac34$ en $e^t = 3$, c'est-à-dire $t = \ln 3$,
donnant
\[
\P\Bigl(S_n \geq \frac{3n}4\Bigr)
\leq \Bigl(\frac{4}{2}\Bigr)^n 3^{-3n/4}
= \bigl(2 \cdot 3^{-3/4}\bigr)^n \approx (0.877)^n ,
\]
exponentiellement petite. La hiérarchie Markov $\to$ Tchebychev $\to$
Chernoff est l'échelle standard~: chaque barreau applique Markov à une
fonction de la variable à croissance plus rapide.
\end{solution}

\begin{solution}{exo:b2:randomvar:8}
Par le théorème de transfert (\cref{thm:b2:randomvar:transfer})
appliqué à $f\bigl(\frac{S_n}{n}\bigr)$ avec $S_n \sim
\mathcal{B}(n, x)$~:
\[
\E\Bigl[f\Bigl(\frac{S_n}{n}\Bigr)\Bigr]
= \sum_{k=0}^n f\Bigl(\frac kn\Bigr)\binom nk x^k(1-x)^{n-k}
= B_nf(x) .
\]
Fixons $\delta > 0$ et découpons $\abs{f(S_n/n) - f(x)}$ selon l'\omterm{def:b2:proba:space}{événement} $D
= \bigl\{\abs{\frac{S_n}{n} - x} \geq \delta\bigr\}$~:
en dehors de $D$, la différence est au plus le module de continuité
$\omega_f(\delta) = \sup_{\abs{s - t}\leq\delta}\abs{f(s) -
f(t)}$~; sur $D$, au plus $2\norm f_\infty$. En prenant les \omterm{def:b2:randomvar:expectation}{espérances} et
en utilisant Tchebychev avec $V\bigl(\frac{S_n}{n}\bigr) =
\frac{x(1-x)}{n} \leq \frac{1}{4n}$~:
\[
\abs{B_nf(x) - f(x)}
\leq \E\,\abs{f(S_n/n) - f(x)}
\leq \omega_f(\delta)
+ 2\norm f_\infty\,\P(D)
\leq \omega_f(\delta) + \frac{2\norm f_\infty}{4n\delta^2} .
\]
La continuité uniforme de $f$ sur $[0, 1]$ rend $\omega_f(\delta) \to
0$~: choisir $\delta$ puis $n$, et $B_nf \to f$ \omterm{def:b2:funcseq:def}{uniformément} --- le
théorème d'approximation de Weierstrass du~\cref{ch:b2:funcseq}, dont le
«~lemme de comptage~» est désormais reconnaissable comme l'inégalité de Tchebychev
pour la \omterm{def:b2:randomvar:law}{loi} binomiale.
\end{solution}

\begin{solution}{exo:b2:randomvar:9}
Développons $S_n^4 = \sum_{i,j,k,l}X_iX_jX_kX_l$ et prenons les
\omterm{def:b2:randomvar:expectation}{espérances}. Par indépendance et centrage, tout terme contenant
un indice apparaissant exactement une fois s'annule ($\E(X_i) = 0$ se
factorise). Termes survivants~: les $n$ termes diagonaux $\E(X_i^4)$, et les
termes appariant deux paires d'indices égaux, $\E(X_i^2X_j^2) =
\E(X_1^2)^2$ pour $i \neq j$, apparaissant $3n(n-1)$ fois~: choisir
la paire non ordonnée de valeurs ($\binom n2$ façons), puis les
$\frac{4!}{2!\,2!} = 6$ façons de les placer dans les quatre emplacements ---
$6\binom n2 = 3n(n-1)$. Donc, avec $\E(X_1^2)^2
\leq \E(X_1^4)$ (Jensen ou Cauchy--Schwarz),
\[
\E(S_n^4) = n\,\E(X_1^4) + 3n(n-1)\,\E(X_1^2)^2
\leq C n^2,
\qquad C = 4\,\E(X_1^4) .
\]
Markov à l'\omterm{def:b2:structures:generated}{ordre} 4~:
\[
\P\Bigl(\Bigl|\frac{S_n}{n}\Bigr| \geq \varepsilon\Bigr)
= \P\bigl(S_n^4 \geq n^4\varepsilon^4\bigr)
\leq \frac{Cn^2}{n^4\varepsilon^4}
= \frac{C}{n^2\varepsilon^4} ,
\]
une série \omterm{def:b2:series:summable}{sommable}. Par Borel--Cantelli~1
(\cref{thm:b2:proba:borelcantelli}), pour chaque $j$ l'\omterm{def:b2:proba:space}{événement}
$B_j = \limsup_n\bigl\{\abs{S_n/n} \geq \frac1j\bigr\}$ a pour
probabilité $0$, donc $\P\bigl(\bigcup_j B_j\bigr) = 0$ par
sous-additivité \omterm{def:b2:structures:countable}{dénombrable}. Sur le complémentaire --- de probabilité $1$
--- pour tout $j$ il existe $N$ tel que $\abs{S_n/n} < \frac1j$ pour
tout $n \geq N$~: précisément $S_n/n \to 0$. La \omterm{def:b2:randomvar:law}{loi} forte des grands
nombres est vraie sous un moment d'\omterm{def:b2:structures:generated}{ordre} quatre~; en retirer cette hypothèse
(théorème de Kolmogorov) relève de la troisième année.
\end{solution}

\begin{solution}{exo:b2:randomvar:10}
$\P(M \leq k) = \bigl(\frac k6\bigr)^2$ (les deux dés valent au plus
$k$, indépendamment), donc $\P(M \geq k) = 1 -
\bigl(\frac{k-1}6\bigr)^2$ et
\[
\E(M) = \sum_{k=1}^6\P(M \geq k)
= 6 - \frac{0 + 1 + 4 + 9 + 16 + 25}{36}
= 6 - \frac{55}{36} = \frac{161}{36} \approx 4.47 ,
\]
confortablement au-dessus de la moyenne $3.5$ d'un seul dé, comme un
maximum se doit de l'être.
\end{solution}

\begin{solution}{exo:b2:randomvar:11}
Avec $I_i = \mathbf 1_{\sigma(i) = i}$~: $\P(\sigma(i) = i) =
\frac{(n-1)!}{n!} = \frac1n$, donc $\E(F_n) = n\cdot\frac1n =
1$. Pour $i \neq j$~: $\P(\sigma(i) = i, \sigma(j) = j) =
\frac{(n-2)!}{n!} = \frac1{n(n-1)}$, d'où
\[
\operatorname{Cov}(I_i, I_j) = \frac1{n(n-1)} - \frac1{n^2}
= \frac{1}{n^2(n-1)} .
\]
Par la boîte à outils de la \omterm{def:b2:randomvar:variance}{variance}
(\cref{thm:b2:randomvar:variancerules}),
\[
V(F_n) = n\cdot\frac1n\Bigl(1 - \frac1n\Bigr)
+ n(n-1)\cdot\frac1{n^2(n-1)}
= 1 - \frac1n + \frac1n = 1 .
\]
Moyenne $1$, \omterm{def:b2:randomvar:variance}{variance} $1$, indépendantes de $n$ --- cohérent
avec la limite de Poisson du problème des rencontres
(\cref{exo:b2:proba:5}).
\end{solution}

\begin{solution}{exo:b2:randomvar:12}
$T_n = \sum_{k=1}^nG_k$ où $G_k \sim \mathcal G(k/n)$ est
le temps pour voir un nouveau jouet quand $k$ manquent, les étapes
étant indépendantes. Donc
\[
V(T_n) = \sum_{k=1}^n\frac{1 - k/n}{(k/n)^2}
\leq \sum_{k=1}^n\frac{n^2}{k^2}
\leq \frac{\pi^2}6\,n^2 ,
\]
par l'\cref{ex:b2:fourier:basel}. Avec $\E(T_n) = nH_n$, $H_n =
\sum_1^n\frac1k$ (\cref{exo:b2:randomvar:3}), Tchebychev donne,
pour $\varepsilon > 0$,
\[
\P\bigl(\abs{T_n - nH_n} \geq \varepsilon\,n\ln n\bigr)
\leq \frac{\pi^2n^2/6}{\varepsilon^2n^2\ln^2 n}
= \frac{\pi^2}{6\,\varepsilon^2\ln^2n}
\xrightarrow[n\to\infty]{} 0 .
\]
Puisque $H_n \sim \ln n$, diviser par $n\ln n$ montre que
$T_n/(n\ln n) \to 1$ en probabilité~: les fluctuations de
$T_n$ sont d'\omterm{def:b2:structures:generated}{ordre} $n$, négligeables devant la moyenne $n\ln n$.
\end{solution}

\begin{solution}{pb:b2:randomvar:1}
\textbf{1.} $\E(S_n) = \frac n2$ et Markov
(\cref{thm:b2:randomvar:markov}) donnent $\P(S_n \geq an) \leq
\frac{n/2}{an} = \frac1{2a}$. Markov ne connaît que la moyenne~: il
ne peut distinguer une variable concentrée en $n/2$ d'une variable
étalée entre $0$ et $n$, de sorte qu'il tarife la queue comme si toute
la masse pouvait s'y trouver.

\textbf{2.} La binomiale équilibrée est \omterm{def:b2:quadratic:adjoint}{symétrique} autour de $n/2$
($S_n$ et $n - S_n$ ont la même \omterm{def:b2:randomvar:law}{loi}), donc avec $x = n(a -
\frac12) > 0$ les deux \omterm{def:b2:proba:space}{événements} $\{S_n - \frac n2 \geq x\}$ et
$\{S_n - \frac n2 \leq -x\}$ sont disjoints et équiprobables~:
$\P(S_n \geq an) = \frac12\P(\abs{S_n - \frac n2} \geq x)$.
Tchebychev avec $V(S_n) = \frac n4$~:
\[
\P(S_n \geq an)
\leq \frac12\cdot\frac{n/4}{n^2(a - 1/2)^2}
= \frac1{8n(a - 1/2)^2},
\]
soit $\frac2n$ en $a = \frac34$.

\textbf{3.} Par indépendance et le théorème du produit,
$\E(\eu^{tS_n}) = \bigl(\E \eu^{tX_1}\bigr)^n = \bigl(\frac{1
+ \eu^t}2\bigr)^n$. Markov appliquée à $\eu^{tS_n}$~:
\[
\P(S_n \geq an) \leq \eu^{-tan}\Bigl(\frac{1 +
\eu^t}2\Bigr)^{\!n} = \exp\Bigl(n\bigl(\ln\tfrac{1 +
\eu^t}2 - ta\bigr)\Bigr).
\]
La dérivée de l'exposant en $t$ est $\frac{\eu^t}{1 + \eu^t}
- a$, s'annulant en $\eu^t = \frac a{1-a}$, c'est-à-dire $t^* =
\ln\frac a{1-a} > 0$~; là $\frac{1 + \eu^{t^*}}2 =
\frac1{2(1-a)}$ et l'exposant vaut
\[
n\Bigl(-\ln 2 - \ln(1-a) - a\ln\frac a{1-a}\Bigr)
= -n\bigl(\ln2 + a\ln a + (1-a)\ln(1-a)\bigr) = -n\,I(a),
\]
avec $I(\frac12) = 0$ et $I'(a) = \ln\frac a{1-a} > 0$ sur
$\intoo{\frac12}1$~: $I(a) > 0$. En $a = \frac34$~:
$\eu^{-I(3/4)} = \frac12(\tfrac34)^{-3/4}(\tfrac14)^{-1/4} =
2\cdot3^{-3/4}$, la borne de l'\cref{exo:b2:randomvar:7}.

\textbf{4.} Les $n + 1$ nombres $\binom nja^j(1-a)^{n-j}$ somment
à $1$, et le plus grand est celui en $j = k = an$ (le mode
de $\mathcal B(n, a)$ est $\floor{(n+1)a} = k$ ici). Un
maximum de $n + 1$ nombres sommant à $1$ vaut au moins
$\frac1{n+1}$~:
\[
\binom nk a^k(1-a)^{n-k} \geq \frac1{n+1}
\quad\Longrightarrow\quad
\binom nk \geq \frac{a^{-an}(1-a)^{-n(1-a)}}{n+1}
= \frac{\eu^{nH(a)}}{n+1}.
\]
Donc $\P(S_n \geq an) \geq \binom{n}{an}2^{-n} \geq
\eu^{n(H(a) - \ln2)}/(n+1) = \eu^{-nI(a)}/(n+1)$~: à un
facteur polynomial $n + 1$ près, l'exposant de Chernoff est la vérité.

\textbf{5.} $n = 100$, $a = \frac34$~: Markov $\frac23$~;
Tchebychev $\frac2{100} = 0.02$~; Chernoff
$(2\cdot3^{-3/4})^{100} = \eu^{-100\,I(3/4)} \approx
2.1\cdot10^{-6}$, contre l'exact $2.8\cdot10^{-7}$. Morale~:
chaque moment d'information divise la borne polynomialement~;
le moment exponentiel change sa \emph{nature}.

\textbf{6.} $\cosh t = \sum_{k\geq0}\frac{t^{2k}}{(2k)!}$ et
$\eu^{t^2/2} = \sum_{k\geq0}\frac{t^{2k}}{2^kk!}$~; la
propriété résulte terme à terme de $(2k)! \geq 2^kk!$, qui est vraie par
récurrence~: $(2k)! = 2k(2k-1)\cdot(2k-2)! \geq 2k\cdot
2^{k-1}(k-1)! = 2^kk!\cdot(2k-1) \geq 2^kk!$.

\textbf{7.} Par indépendance, $\E\bigl(\eu^{t\sum\varepsilon_i}
\bigr) = (\cosh t)^n \leq \eu^{nt^2/2}$, donc Markov donne
$\P(\sum\varepsilon_i \geq s) \leq \eu^{nt^2/2 - ts}$~;
en minimisant en $t = s/n$ on obtient $\eu^{-s^2/(2n)}$.

\textbf{8.} Avec $X_i = \frac{1 + \varepsilon_i}2$, $\widehat
p_n - \frac12 = \frac1{2n}\sum\varepsilon_i$, donc
$\{\widehat p_n - \frac12 \geq \delta\} =
\{\sum\varepsilon_i \geq 2n\delta\}$ et la question 7 donne la
borne $\eu^{-(2n\delta)^2/(2n)} = \eu^{-2n\delta^2}$. L'\omterm{def:b2:proba:space}{événement}
\omterm{def:b2:quadratic:adjoint}{symétrique} a la même borne, d'où le facteur $2$ pour
$\abs{\widehat p_n - \frac12} \geq \delta$.

\textbf{9.} $\E(\eu^{tX}) = \sum_x\eu^{tx}\P(X = x)$ est une
série de fonctions lisses de $t$ dont les dérivées terme à terme
sont dominées, sur tout intervalle \omterm{def:b2:metric:compact}{compact} en $t$, par
$\eu^{\abs t}\P(X = x)$ (car $0 \leq x \leq 1$)~: par le
théorème de dérivation pour les séries \omterm{def:b2:funcseq:series}{normalement} convergentes
(\cref{thm:b2:funcseq:differentiation}) elle est deux fois
dérivable, et la règle du quotient donne $\psi' = \E_t(X)$
et $\psi'' = \E_t(X^2) - \E_t(X)^2$, où $\E_t$ est
l'\omterm{def:b2:randomvar:expectation}{espérance} pour les poids repondérés
$\eu^{tx}\P(X{=}x)/\E(\eu^{tX})$ --- positifs, sommant à
$1$, portés par les mêmes valeurs $x \in \intcc01$. Une \omterm{def:b2:randomvar:variance}{variance}
d'une variable à valeurs dans $\intcc01$ est au plus $\frac14$~: par
l'\cref{exo:b2:randomvar:6}, elle vaut $\min_c\E_t((X - c)^2) \leq
\E_t\bigl((X - \tfrac12)^2\bigr) \leq \tfrac14$. Taylor avec
reste intégral, en utilisant $\psi(0) = 0$, $\psi'(0) = p$~:
\[
\psi(t) = tp + \int_0^t(t - s)\,\psi''(s)\,\dd s
\leq tp + \frac{t^2}2\cdot\frac14,
\]
c'est-à-dire $\E(\eu^{t(X - p)}) \leq \eu^{t^2/8}$ pour tout $t$ réel.

\textbf{10.} Par indépendance, $\E\bigl(\eu^{t(S_n -
np)}\bigr) \leq \eu^{nt^2/8}$~; Markov et l'optimisation $t
= 4\delta$ donnent
\[
\P(\widehat p_n - p \geq \delta)
\leq \eu^{nt^2/8 - tn\delta}\Big|_{t = 4\delta}
= \eu^{-2n\delta^2};
\]
en appliquant ceci aux variables $1 - X_i$ (également dans
$\intcc01$) on borne l'autre queue, d'où le
$2\eu^{-2n\delta^2}$ bilatéral.

\textbf{11.} Tchebychev ne requiert qu'un moment d'\omterm{def:b2:structures:generated}{ordre} deux et donne
$\frac{p(1-p)}{n\delta^2}$~; Hoeffding requiert le
\emph{caractère borné} et donne $2\eu^{-2n\delta^2}$. En $p =
\frac12$, $\delta = 0.03$~: les bornes sont
$\frac{278}{n}$ (approximativement) contre
$2\eu^{-0.0018n}$~; elles se croisent vers $n \approx 1200$, après
quoi la borne exponentielle l'emporte, et démesurément ($n = 5000$~:
$0.056$ contre $2.5\cdot10^{-4}$).

\textbf{12.} Par Hoeffding (question 10), $\P(\abs{\widehat
p_n - p} \geq \delta) \leq 2\eu^{-2n\delta^2} \leq \alpha$ dès
que $2n\delta^2 \geq \ln\frac2\alpha$, c'est-à-dire $n \geq
\frac{\ln(2/\alpha)}{2\delta^2}$.

\textbf{13.} $\delta = 0.03$, $\alpha = 0.05$~: $n \geq
\frac{\ln 40}{2\cdot0.0009} \approx 2049.4$~: $2050$ personnes.
Pour $\delta = 0.01$~: $n \geq \frac{\ln40}{0.0002} \approx
18\,445$. La taille de la population n'apparaît jamais parce que chaque
électeur échantillonné est modélisé comme un nouveau tirage Bernoulli$(p)$~: la
difficulté du sondage est la \omterm{def:b2:randomvar:variance}{variance} d'une pièce, non la taille
du pays. Diviser la marge par deux coûte quatre fois l'échantillon
--- la \omterm{def:b2:randomvar:law}{loi} en $1/\delta^2$.

\textbf{14.} Tchebychev~: $\P(\abs{\widehat p_n - p} \geq
\delta) \leq \frac{p(1-p)}{n\delta^2} \leq
\frac1{4n\delta^2} \leq \alpha$ pour $n \geq
\frac1{4\alpha\delta^2}$, c'est-à-dire $5556$ à trois points ---
environ $2.7$ fois l'exigence de Hoeffding. Sans
remise, la \omterm{def:b2:randomvar:variance}{variance} est multipliée par $\frac{N -
n}{N-1} < 1$ (\cref{exo:b2:randomvar:5}), donc le même $n$ ne
peut que faire mieux~: le calcul avec remise est le
conservateur.

\textbf{15.} $\{\widehat p_n \leq \frac12\} \subseteq
\{\widehat p_n - 0.52 \leq -0.02\}$, donc par la borne de Hoeffding
unilatérale $\P(\widehat p_n \leq \tfrac12) \leq
\eu^{-2n(0.02)^2} \leq 0.01$ dès que $n \geq \frac{\ln
100}{2\cdot0.0004} \approx 5756.5$~: $5757$ électeurs. Le coût
varie comme l'inverse du carré de l'\emph{avance}, non de la
précision désirée~: les courses serrées sont chères.

\textbf{16.} Utilisé~: l'échantillon est tiré \omterm{def:b2:funcseq:def}{uniformément} et
indépendamment de l'électorat~; toute personne échantillonnée
répond, honnêtement, et $p$ ne bouge pas pendant le sondage. Les vrais
sondages violent les trois~: les répondants joignables et disposés
ne sont pas un échantillon uniforme (biais de sélection et de non-réponse),
et les réponses peuvent être mensongères ou instables. Ce sont des
erreurs de \emph{biais}~: elles décalent $\E(\widehat p_n)$ loin de
$p$ d'une quantité indépendante de $n$, de sorte qu'aucune taille
d'échantillon ne les réduit --- les mathématiques de cette partie ne contrôlent que
le terme de fluctuation.

\textbf{17.} Tchebychev pour un groupe de taille $m$~: $\P(\abs{
\widehat p^{(i)} - p} \geq \delta) \leq \frac{1}{4m\delta^2}
\leq \frac18$ pour $m \geq \frac2{\delta^2}$. Si moins de
$k/2$ groupes se trompent, alors plus de $k/2$ des valeurs
$\widehat p^{(i)}$ se situent dans l'intervalle \omterm{def:b2:metric:topology}{ouvert} $\intoo{p -
\delta}{p + \delta}$, et leur médiane aussi~; donc
$\{\abs{M - p} \geq \delta\}$ force au moins $\lceil
k/2\rceil$ erreurs parmi $k$ groupes \omterm{def:b2:proba:independence}{indépendants}. La borne de l'union
sur les $\binom k{\lceil k/2\rceil}$ ensembles possibles de
groupes fautifs donne
\[
\P(\abs{M - p} \geq \delta)
\leq \binom{k}{\lceil k/2\rceil}
\Bigl(\frac18\Bigr)^{k/2}
\leq 2^k\,8^{-k/2} = 2^{-k/2} :
\]
décroissance exponentielle en le nombre de groupes, achetée avec
rien d'autre que des \omterm{def:b2:randomvar:variance}{variances} --- utile précisément quand les termes
sont non bornés et que Hoeffding est indisponible.

\textbf{18.} Cauchy--Schwarz
(\cref{thm:b2:randomvar:jensen})~:
\[
\E(X) = \E(X\,\mathbf 1_{X>0})
\leq \sqrt{\E(X^2)}\sqrt{\E(\mathbf 1_{X>0}^2)}
= \sqrt{\E(X^2)\,\P(X > 0)} ;
\]
élever au carré et diviser.

\textbf{19.} Soit $h(a) = I(a) - 2(a - \tfrac12)^2$. Alors
$h(\tfrac12) = 0$, $h'(a) = \ln\frac a{1-a} - 4(a -
\tfrac12)$ s'annule en $\tfrac12$, et
\[
h''(a) = \frac1a + \frac1{1-a} - 4 = \frac{1}{a(1-a)} - 4
\geq 0
\]
puisque $a(1-a) \leq \frac14$. Donc $h'$ croît depuis $0$ sur
$\intco{\frac12}1$, d'où $h' \geq 0$ et $h \geq 0$~: $I(a)
\geq 2(a - \tfrac12)^2$.

\textbf{20.} $I(\tfrac12) = I'(\tfrac12) = 0$, $I''(a) =
\frac1{a(1-a)}$ donne $I''(\tfrac12) = 4$, et $I'''(\tfrac12)
= 0$ (la fonction est \omterm{def:b2:quadratic:adjoint}{symétrique} autour de $\tfrac12$), donc
$I(\tfrac12 + \delta) = 2\delta^2 + O(\delta^4)$. La question 4
borne alors la vraie queue \emph{en dessous} par
$\eu^{-n(2\delta^2 + O(\delta^4))}/(n+1)$~: pour $\delta$ petit
l'exposant de Hoeffding $2n\delta^2$ est asymptotiquement exact
--- seules des améliorations polynomiales en $n$ sont possibles.

\textbf{21.} Markov~: un moment, décroissance $1/a$, utile seulement comme
moteur derrière les autres (la question 1 le montre plat).
Tchebychev~: deux moments, décroissance $\frac{V}{n\delta^2}$, optimale
sans hypothèses supplémentaires
(\cref{ex:b2:randomvar:chebsharp}), et le meilleur outil à la
question 23. Moment d'\omterm{def:b2:structures:generated}{ordre} quatre (\cref{exo:b2:randomvar:9})~:
décroissance $C/n^2$, juste assez de sommabilité pour une \omterm{def:b2:randomvar:law}{loi} forte.
Hoeffding~: variables bornées, décroissance $2\eu^{-2n\delta^2}$,
le cheval de trait de la partie III. Chernoff avec le taux exact
$I(a)$~: moments exponentiels \omterm{def:b2:metric:complete}{complets}, exposant imbattable
(questions 4, 20), le point de référence pour tout le reste.

\textbf{22.} Si la pièce est équilibrée~: $\P(\widehat p_n > 0.525)
\leq \P(\widehat p_n - \tfrac12 \geq 0.025) \leq
\eu^{-2n(0.025)^2}$. Si $p = 0.55$~: $\P(\widehat p_n \leq
0.525) \leq \P(\widehat p_n - 0.55 \leq -0.025) \leq
\eu^{-2n(0.025)^2}$. Les deux erreurs sont sous $0.01$ lorsque
$2n(0.025)^2 \geq \ln 100$, c'est-à-dire $n \geq 3684.2$~: $3685$
lancers. (Distinguer des hypothèses distantes de $2.5$ points coûte
ce qu'estimer à $\pm2.5$ points coûte.)

\textbf{23.} Hoeffding~: $n \geq \frac{\ln 40}{2(0.005)^2}
\approx 73\,778$. Tchebychev avec la vraie \omterm{def:b2:randomvar:variance}{variance} $p(1-p)
= 0.0099$~: $n \geq \frac{0.0099}{0.05\cdot(0.005)^2} =
7920$ --- neuf fois moins cher. L'exposant de Hoeffding
$2n\delta^2$ tarife la \omterm{def:b2:randomvar:variance}{variance} à son pire cas $\frac14$,
absurdement pessimiste quand $p = 0.01$~; le humble moment
d'\omterm{def:b2:structures:generated}{ordre} deux sait mieux. L'outil manquant est une borne exponentielle
consciente de la \omterm{def:b2:randomvar:variance}{variance} (inégalité de Bernstein, troisième année) --- ou,
pour les \omterm{def:b2:proba:space}{événements} rares, l'approximation de Poisson démontrée au
\cref{ch:b2:genfun}, qui opère sur l'échelle relative naturelle.

\textbf{24.} Fixons $\delta > 0$~: $\sum_n 2\eu^{-2n\delta^2} <
\infty$ (série de type géométrique), donc Borel--Cantelli~1
(\cref{thm:b2:proba:borelcantelli}) donne
$\P(\abs{\widehat p_n - p} \geq \delta \text{ infiniment
souvent}) = 0$, c'est-à-dire l'\omterm{def:b2:proba:space}{événement} $E_j =
\bigcup_N\bigcap_{n\geq N}\{\abs{\widehat p_n - p} <
\tfrac1j\}$ a pour probabilité $1$ pour chaque $j$. L'intersection
\omterm{def:b2:structures:countable}{dénombrable} $\bigcap_jE_j$ a encore probabilité $1$
(sous-additivité sur les complémentaires), et dessus $\widehat p_n
\to p$~: la \omterm{def:b2:randomvar:law}{loi} forte des grands nombres pour les lancers de pièce, le
caractère borné jouant le rôle que le moment d'\omterm{def:b2:structures:generated}{ordre} quatre a joué à
l'\cref{exo:b2:randomvar:9}.

\textbf{25.} Un moment achète une borne plate~; deux achètent
$1/(n\delta^2)$, et rien de plus (l'exemple d'optimalité)~; quatre
achètent $1/n^2$, assez pour se télescoper en une \omterm{def:b2:randomvar:law}{loi} presque sûre~;
le caractère borné achète $\eu^{-2n\delta^2}$~; et le moment
exponentiel \omterm{def:b2:metric:complete}{complet} achète le taux exact $I$, qu'aucune méthode
ne bat. Sonder $2050$ personnes suffit pour n'importe quel pays parce que
la fluctuation de l'échantillon est régie par la \omterm{def:b2:randomvar:variance}{variance} de la pièce,
non par la taille de la population --- les étiquettes de prix en $1/\delta^2$ et
$\ln(1/\alpha)$ sont universelles. Le théorème central limite du
volume de troisième année remplace ces inégalités, sur l'échelle
$\sqrt n$, par une \omterm{def:b2:randomvar:law}{loi} limite \omterm{def:b2:multint:exact}{exacte} à constantes explicites ---
transformant chaque borne de ce problème en une égalité asymptotique.
\end{solution}
